\subsection{Выборка, вход модели и извлечение программы}
Исходный набор BigCodeBench содержит 1\,140 задач. Идентификаторы сортируются как строки и перемешиваются генератором Python \texttt{random.Random(20260908)}. Первые 40 образуют выборку разработки. Из оставшейся последовательности основное исследование берёт первые 200 идентификаторов после исключения /0--/7; опыт с оформлением инструкций использует следующие 80 с теми же исключениями. Состав выборок указан в полных таблицах исходов. В основном опыте каждая из четырёх групп по 50 задач сначала сортируется по идентификатору; затем с тем же начальным состоянием генератора перемешиваются блоки задач и пять условий внутри блока. Метки повторов различают новые генерации и не управляют случайным состоянием самой модели.

Вход BigCodeBench состоит из исходного поля \texttt{instruct\_prompt} в блоке задачи и поля \texttt{code\_prompt}, включая стартовый код, в блоке предоставленного контекста. Поля \texttt{test} и \texttt{canonical\_solution} модели не передаются. Приложение с запросами задаёт инструкции этапов, обозначения, порядок блоков и требуемый формат ответа. В опыте с оформлением инструкция условия помещается после текста задачи и перед блоком контекста.

Для D, SN, SR, MN, MR и INoT* извлечение распознаёт блоки, ограждённые тройными обратными кавычками и помеченные \texttt{python} или \texttt{py} без учёта регистра. После открывающей и перед закрывающей границей должен стоять перевод строки. Требуется ровно один распознанный блок; блоки с другими языковыми метками оценщик извлечения игнорирует. Начальные и конечные пробельные символы удаляются. Если подходящих блоков нет или их несколько, создаётся пустая программа с признаком ошибки формата. Произвольное оформление ответа не интерпретируется и не исправляется. При наличии завершённого ответа пустая программа проходит оценку и учитывается как неуспех; отсутствующая генерация остаётся пропуском. В SCC используется авторская процедура извлечения, описанная ниже.

\subsection{Штатные тесты и окружение}
Основная оценка использует Python 3.11, задачи BigCodeBench v0.1.4, локальное выполнение, режим \texttt{instruct full}, два параллельно работающих процесса оценщика и выбранные идентификаторы задач. Автоматическая вставка стартового кода отключена (\texttt{calibrated=False}); автоматическую генерацию эталонов заменяют отдельные предварительные контроли (\texttt{no\_gt=True}). Минимальное время задано как 0,1 секунды, однако фактический предел оценщика для набора тестов одного кандидата составляет 241 секунду. Он отличается от 180-секундного ограничения вызова модели. Пределы процесса составляют 30\,720 МБ для адресного пространства и данных и 10 МБ для стека; контейнер дополнительно ограничен 3 ГБ суммарной памяти, двумя CPU и 256 процессами. Временная файловая система имеет объём 512 МБ.

Основное окружение использует NumPy 1.26.4, pandas 2.2.3, SciPy 1.14.1, scikit-learn 1.5.2, Matplotlib 3.9.2 и NLTK 3.8.1. На этапе контроля добавлены BeautifulSoup 4.12.3, OpenCV-headless 4.9.0.80, mechanize 0.4.10, rsa 4.9, pyquery 2.0.0 и statsmodels 0.14.1. Ресурсы NLTK \texttt{punkt}, \texttt{averaged\_perceptron\_tagger}, \texttt{stopwords} и \texttt{words} установлены до тестирования без сети. Отдельное окружение опыта на 80 задачах также обеспечивает совместимость со старыми библиотечными интерфейсами, включая SciPy 1.10.1, scikit-learn 1.3.1 и Matplotlib 3.7.0, и содержит TensorFlow-CPU 2.15.1. Поэтому результаты этого опыта анализируются отдельно.

Эталонную программу получают объединением \texttt{code\_prompt} и \texttt{canonical\_solution}. Для отрицательного контроля в её конец добавляется замена \texttt{task\_func(*args, **kwargs)}, возвращающая \texttt{None}. Задача допускается к анализу качества, только если эталон проходит, а заведомо неверная программа не проходит тесты. Такой отрицательный контроль оказался нечувствительным на /59 в опыте с 80 задачами: успешного выполнения эталона здесь недостаточно для включения в анализ качества. На допустимых задачах штатный успех кодируется как 1, штатный отказ или превышение времени тестирования --- как 0. Отсутствующая генерация, незавершённый запуск оценщика или непригодный контроль не получают вымышленного бинарного исхода.

\subsection{Расчёты по парным исходам}
Для завершённого сравнения с одним повтором обозначим через $b$ число задач, решённых только условием A, через $c$ --- только условием B, через $n$ --- число всех полных допустимых пар. Разность качества равна $(b-c)/n$. При $d=b+c>0$ точное двустороннее значение теста Мак-Немара имеет вид
\[
p=\min\left(1,\;2\sum_{k=0}^{\min(b,c)}{d\choose k}2^{-d}\right);
\]
при $d=0$ принимается $p=1$. В семействе из $m$ тестов исходные значения упорядочиваются как $p_{(1)}\le\cdots\le p_{(m)}$, после чего вычисляется
\[
p^{\mathrm{Holm}}_{(i)}
=\min\left(1,\max_{j\le i}(m-j+1)p_{(j)}\right).
\]
В основном опыте $m=4$, в опыте с оформлением $m=2$. Границы 95\%-интервала --- процентили 2,5 и 97,5 распределения средних в 10\,000 повторных выборках задач, полученных генератором NumPy с начальным состоянием 20260908. Все условия выбранной задачи переносятся в повторную выборку вместе. Интервалы служат описанием неопределённости; заранее заданные решения об отклонении гипотез принимаются по тестам с поправкой Холма.

Если среди $N$ назначений наблюдаются $s$ успехов и $u$ недоступных исходов, крайние границы доли успехов равны $[s/N,(s+u)/N]$. Для парной разности каждый неизвестный бинарный исход отдельно заменяется на 0 или 1 так, чтобы получить минимум и максимум возможной разности. Эти границы характеризуют неопределённость пропусков, а не случайность выборки.

\subsection{Адаптация SCC и обратная связь}
Процесс SCC в реализации 2024 года начинается с краткого плана аналитика в формате JSON, по которому разработчик создаёт программу Python. Тестировщик генерирует функцию \texttt{check(candidate)} с печатью результатов не более пяти примеров; инструкция явно запрещает проверки \texttt{assert}. Контроллер выполняет программу вместе со сгенерированной проверкой. Отсутствие исключений даёт внутреннее сообщение об успехе; исключение или превышение времени формирует обратную связь для разработчика. Напечатанные ответы автоматически не сравниваются с ожидаемыми. Внутренний предел выполнения составляет 10 секунд; на процесс контейнера отдельно отводится 45 секунд. Эталонные решения и тесты бенчмарка в эту обратную связь не входят.

Предел в две итерации разработчика допускает начальную программу и не более одного исправления: аналитик, разработчик, тестировщик и, при необходимости, снова разработчик. Последний разрешённый ответ разработчика завершает сессию без ещё одного вызова тестировщика. Авторская процедура извлекает первый ограждённый блок кода, а при его отсутствии пытается восстановить код по определениям функций. Контроллер сохраняет последнюю программу с распознаваемой функцией верхнего уровня либо ошибку, если такой программы нет. История сообщений каждой роли целиком передаётся модели как JSON-диалог. Исходные параметры API для температуры, top-$p$ и максимального числа выходных токенов подписочным способом вызова не поддерживаются. Пилот оценивает именно эту адаптацию; внутреннее сообщение SCC об успехе не является отдельно измеряемым результатом BigCodeBench.

\subsection{Факторный эксперимент с повторами и дополнительная серия SCC}
\label{app:expanded-studies}
Факторная серия Luna включает 200 задач первого эксперимента и ещё 800 задач из исходного случайного порядка, без разработческой выборки, 80 задач исследования разметки и /0--/7. Ранее изученные 200 задач не являются новой отложенной выборкой. Все запросы выполняются заново; прежние программы не добавляются к новым наблюдениям. Из 15\,000 назначений получена 14\,961 завершённая программа. Контрольные проверки до генерации разрешают измерять качество на 985 задачах. Все назначения, включая 15 задач с неуспешным контролем среды, остаются в учёте.

\begin{table}[htbp]\centering\small
\caption{Выборки Luna. Блок объединяет задачу и метку повтора; в нём назначены все методы. Завершённая программа не обязательно проходит тесты.}
\begin{tabularx}{\textwidth}{@{}lXX@{}}\toprule
 & Факторная серия & Дополнительная серия SCC\\\midrule
Идентификаторы задач & 1\,000 & 956 из той же выборки\\
Повторы на задачу & Три & От одной до трёх выбранных меток\\
Условия & D, SN, SR, MN, MR & Новые SR, SN и SCC\\
Назначения & 15\,000 & 6\,000 в 2\,000 блоках\\
Максимум вызовов & 27\,000 & 12\,000\\
Семейство тестов качества & Четыре сравнения & SCC--SR и SCC--SN\\\bottomrule
\end{tabularx}\end{table}

Обе серии используют GPT-5.6 Luna, уровень рассуждения medium, Codex CLI 0.153.4, срок 600 секунд на вызов и не более восьми параллельных исполнителей. Запрос длиной более 65\,536 байт UTF-8 отклоняется, а не сокращается. Весь предоставленный контекст сохраняется. Начатые попытки не повторяются; после прерывания разрешено запустить только ещё не начатое назначение. Завершённый ответ неверного формата даёт наблюдаемого пустого кандидата; отсутствие завершённого ответа остаётся пропуском. Имя модели и повторные запросы не доказывают неизменность весов или независимость выборок поставщика.

Для входных токенов $I$, кешированной части $C\subseteq I$ и выходных токенов $O$ денежные оценки Luna равны
\[
V_L=\frac{0.20(I-C)+0.02C+1.20O}{10^6},\qquad
V_{L,0}=\frac{0.20I+1.20O}{10^6}\quad\text{долл. США}.
\]
Стандартные тарифы API проверены 13 сентября 2026 года~\cite{openaipricing2026}; они отличаются от тарифов mini в первом эксперименте. Это оценка измеренных счётчиков подписочного канала по опубликованным ставкам, а не счёт за подписку или полные расходы исследования. Отдельно показанные токены рассуждения входят в выход и повторно не прибавляются. Неизвестный расход не заменяется нулём.

Исходный план SCC содержит 3\,000 случайно упорядоченных блоков «задача--повтор» и 9\,000 назначений методов. Из-за ограничения бюджета до просмотра качества SCC зафиксированы первые 2\,000 блоков этого порядка. В них входят 956 задач: 212 с одной выбранной меткой повтора, 444 с двумя и 300 с тремя. Контроль среды допускает оценку качества на 941 из этих задач. Оставшиеся 3\,000 назначений не выбраны по организационным причинам, не входят в знаменатель 6\,000 и сохраняются в исходном реестре. Восемь ранее начатых и прерванных попыток остаются недоступными; новые ответы для них не генерируются. Изменение принято после начала генерации, поэтому оно не описывается как полностью предзарегистрированный новый эксперимент. Выбранный префикс завершён и может быть представлен в результатах: для SCC завершены 1\,600 процессов (396 с ошибкой Docker или инфраструктуры и 4 приостановлены), для SR --- 1\,997 (3 приостановлены), для SN --- 1\,998 (1 с потоковой или инфраструктурной ошибкой). Итого завершены 5\,595 кандидатов; 397 относятся к инфраструктурным ошибкам, 8 --- к приостановленным попыткам. Каждый завершённый кандидат прошёл штатную нативную оценку. Эти числа обозначают завершённые программы или процессы, а не успешные результаты тестов.

В факторной серии задача входит в сравнение качества, только если в обоих условиях известны все три повторных исхода. Разность $d_i$ вычисляется между средними по трём повторам. Двусторонний парный $t$-тест проверяет среднюю разность по задачам. Четыре сравнения таблицы~\ref{tab:scale-contrasts} образуют одно фиксированное семейство Холма при $\alpha=0.05$. Каждый описательный интервал использует 10\,000 выборок целых задач с возвращением, генератор PCG64 из NumPy 1.26.4, seed 20260909 и зафиксированный порядок задач. Границы процентилей определяются линейной интерполяцией. Эти правила отделены от четырёх точных тестов Мак-Немара первого эксперимента с mini.

Дополнительная оценка SCC учитывает неодинаковое число выбранных повторов. Для задачи $i$ и метода сравнения $B$ обозначим через $K_i^B$ выбранные метки повторов, у которых известны оба исхода --- SCC и $B$ --- при действительном контроле среды. Вклад задачи равен
\[
d_i^B=\frac{1}{|K_i^B|}\sum_{r\in K_i^B}
\bigl(Y_{ir}(\mathrm{SCC})-Y_{ir}(B)\bigr).
\]
Учитываются только задачи с $|K_i^B|>0$, и каждая имеет одинаковый вес. Поэтому три выбранных повтора не дают задаче тройного веса относительно задачи с одним повтором. SCC--SR и SCC--SN образуют фиксированное семейство из двух тестов Холма. Интервалы используют 10\,000 выборок целых задач, PCG64, seed 20260911 и линейную интерполяцию процентилей. Отдельно сохраняется исходная оценка SCC, требующая трёх полных повторов в обоих условиях, со своим семейством из двух тестов и исходным бутстрэпом Python Mersenne Twister. Её границы пропусков для 9\,000 назначений не выдаются за границы выбранной серии. Повторы не считаются независимыми задачами; успех лучшего из трёх ответов не оценивается.

Для обеих оценок по задачам статистика равна $t=\bar d/(s_d/\sqrt n)$ при $n\geq2$ и $s_d>0$, с $n-1$ степенями свободы. Это заданное приближение для большой выборки. Если тест SCC нельзя оценить, он сохраняет место в семействе: в поправку передаётся 1, а оснований для отклонения нет. Постоянный нулевой вектор при двух или более задачах даёт $p=1$. Зафиксированная факторная программа кодирует постоянный ненулевой вектор признаком вырождения и $p=0$; это не обычный конечный $t$-тест. В четырёх наблюдаемых факторных сравнениях такого вырождения нет. Ни один анализ не вводит допуск неухудшения качества или совместный критерий качества и денежных затрат.

Границы пропусков и выборочные интервалы отвечают на разные вопросы. Для выбранной серии SCC знаменатель каждой задачи равен числу её \emph{назначенных} выбранных повторов, включая недоступные исходы. Неизвестные бинарные исходы заполняются всеми благоприятными либо неблагоприятными значениями; затем разности усредняются с одинаковым весом задач. Границы приводятся для 941 допустимой задачи и для всех 956 выбранных задач. Они описывают диапазон, разрешённый пропусками, а не доверительный интервал или поправку на неслучайную недоступность. Факторные границы также сохраняют полные назначенные знаменатели.

Для сравнения ресурсов успешный контроль среды не требуется. В факторной серии нужны три завершённых процесса в обоих условиях с полными счётчиками всех вызовов. В дополнительной серии SCC пара выбранных повторов допустима, только если оба процесса целиком завершены и для каждого начатого вызова есть корректные счётчики. Расход суммируется по вызовам, затем усредняется по совпадающим повторам внутри задачи. Задачи вновь получают одинаковый вес. Известный расход незавершённых процессов сообщается отдельно как неполная сумма. Исходная ресурсная оценка SCC по известному расходу начатых вызовов трёх повторов сохраняется отдельно и может учитывать попытки без завершённой программы.

Для парных средних расходов по задачам $A_i$ и $B_i$ различаются два отношения:
\[
R_{\mathrm{pooled}}=\frac{\sum_i A_i}{\sum_i B_i},\qquad
R_{\mathrm{task}}=\frac{1}{n_R}\sum_{i:B_i\ne0}\frac{A_i}{B_i},
\]
где $n_R$ --- число ненулевых знаменателей. Факторная серия использует seed 20260909 для интервалов отношения сумм и средней парной разности. Дополнительный анализ SCC использует seed 20260912 для разности, 20260913 для среднего отношения по задачам и 20260914 для отношения сумм. Для каждой статистики генератор начинается заново; все 10\,000 реализаций сохраняются. Исходный анализ SCC даёт для отношения сумм лишь точечную оценку; его интервал среднего отношения по задачам относится к другой величине. Таблицы называют оценку и число полных пар. Ресурсные интервалы описательны и не входят в семейства тестов качества.

Проверка чувствительности к сходству задач выбирает с возвращением целые группы по четырём заранее заданным разбиениям: совместному правилу инструкций и кода и порогам сходства инструкций 0.50, 0.70 и 0.90. Используется seed 20260911. Сумма разностей выбранных задач делится на число выбранных задач, что сохраняет их вес. Задачи не исключаются по сходству. Эти интервалы дополняют основной анализ по задачам. SCC остаётся сравнением полных процессов, включая обратную связь от его собственных проверок; вариант SCC без обозначений ролей не оценивается.
